\documentclass[10pt,aspectratio=169]{beamer}
\input{../../_en_preambulo_beamer}
% Comandos y macros estadísticas compartidas para las presentaciones Beamer en inglés (Metropolis)

% Conjuntos numéricos básicos
\newcommand{\R}{\mathbb{R}}
\newcommand{\N}{\mathbb{N}}
\newcommand{\Q}{\mathbb{Q}}
\newcommand{\Z}{\mathbb{Z}}
\newcommand{\set}[1]{\left\{ #1 \right\}}
\newcommand{\abs}[1]{\left|#1\right|}
\newcommand{\norm}[1]{\left\| #1 \right\|}

% Comandos estadísticos y probabilísticos del curso
\newcommand{\Var}{\operatorname{Var}}
\newcommand{\cov}{\operatorname{Cov}}
\newcommand{\corr}{\rho}
\newcommand{\comb}[2]{\begin{pmatrix} #1 \\ #2 \end{pmatrix}}
\newcommand{\s}{\sigma}
\newcommand{\particion}{\mathfrak{P}}
\newcommand{\card}[1]{\#\left( #1 \right)}
\newcommand{\seq}[3]{\set{#1}_{#2}^{#3}}
\newcommand{\E}{\mathbb{E}}
\newcommand{\Prob}{\mathbb{P}}

% Entornos pedagógicos y cajas en inglés académico natural (soportando tanto nombres en inglés como en español)
\newenvironment{discussion}{%
  \begin{alertblock}{Discussion Question}%
}{%
  \end{alertblock}%
}
\newenvironment{discusion}{%
  \begin{alertblock}{Discussion Question}%
}{%
  \end{alertblock}%
}

\newenvironment{reflection}{%
  \begin{alertblock}{Classroom Reflection}%
}{%
  \end{alertblock}%
}
\newenvironment{reflexion}{%
  \begin{alertblock}{Classroom Reflection}%
}{%
  \end{alertblock}%
}

\newenvironment{paradox}{%
  \begin{alertblock}{Exploratory Paradox}%
}{%
  \end{alertblock}%
}
\newenvironment{paradoja}{%
  \begin{alertblock}{Exploratory Paradox}%
}{%
  \end{alertblock}%
}

\newenvironment{motivation}{%
  \begin{exampleblock}{Motivation: Data Science in Practice}%
}{%
  \end{exampleblock}%
}
\newenvironment{motivacion}{%
  \begin{exampleblock}{Motivation: Data Science in Practice}%
}{%
  \end{exampleblock}%
}

\newenvironment{quickchallenge}{%
  \begin{block}{Quick Team Challenge}%
}{%
  \end{block}%
}
\newenvironment{retoclase}{%
  \begin{block}{Quick Team Challenge}%
}{%
  \end{block}%
}


\title[Counting Techniques]{Section 02.03: Counting Techniques}
\subtitle{Permutations, combinations, and classical probability}
\author[J. Castillo Colmenares]{Juliho Castillo Colmenares}
\institute{Tecnológico de Monterrey}
\date{\vspace{-1.5cm}}

\begin{document}

% 01 — Identity
\begin{frame}[plain]
  \vspace{-0.8cm}
  \titlepage
\end{frame}

% 02 — Roadmap
\begin{frame}{Roadmap --- Chapter 02: Probability Theory}
  \scriptsize
  \begin{itemize}\itemsep=0.04em
    \item \textbf{02.01} Set theory and probability.
    \item \textbf{02.02} Probability foundations and axioms.
    \item \textbf{02.03} Counting techniques \emph{(today)}.
    \item \textbf{02.04} Conditional probability and Bayes' rule.
    \item \textbf{02.05} Bayes' theorem.
    \item \textbf{02.06} Random sampling.
  \end{itemize}
  \vspace{-0.05cm}
  \begin{block}{Learning objective}
    Count finite sample spaces with the multiplication principle, permutations,
    and combinations, and turn those counts into classical probabilities when
    outcomes are equiprobable.
  \end{block}
\end{frame}

% 03 — Motivation
\begin{frame}{From listing outcomes to systematic counting}
  \small
  \begin{columns}[T]
    \begin{column}{0.48\textwidth}
      \begin{block}{The challenge}
        A sample space can grow too quickly to list every outcome. Counting
        describes it without enumerating it.
      \end{block}
    \end{column}
    \begin{column}{0.48\textwidth}
      \begin{alertblock}{Classical probability}
        If all outcomes are equally likely,
        \[
          P(A)=\frac{\lvert A\rvert}{\lvert\Omega\rvert}.
        \]
        Counting techniques compute both cardinalities.
      \end{alertblock}
    \end{column}
  \end{columns}
\end{frame}

% 04 — Fundamental principle
\begin{frame}{Fundamental counting principle}
  \small
  \begin{block}{Product rule}
    If a first operation has $n_1$ options, a second has $n_2$ options for
    every previous choice, and so on, then
    \[
      N=n_1n_2\cdots n_k.
    \]
  \end{block}
  \pause
  \begin{exampleblock}{Interpretation}
    Each final outcome is a path: multiplication is used because every stage
    combines with all options in later stages.
  \end{exampleblock}
\end{frame}

% 05 — Counting model
\begin{frame}{How to model a counting question}
  \footnotesize
  \begin{enumerate}\itemsep=0.12cm
    \item Describe precisely the outcome being counted.
    \item Separate the process into successive stages or choices.
    \item Ask whether order matters and whether repetition is allowed.
    \item Choose multiplication, permutations, or combinations.
  \end{enumerate}
  \vspace{0.1cm}
  \begin{alertblock}{Control rule}
    If exchanging positions changes the outcome, use an ordered model; if it
    does not, count an unordered selection.
  \end{alertblock}
\end{frame}

% 06 — Permutations
\begin{frame}{Permutations: when order matters}
  \small
  \begin{block}{Definition}
    The number of ordered arrangements of $r$ distinct objects chosen from
    $n$, without repetition, is
    \[
      P(n,r)=\frac{n!}{(n-r)!}=n(n-1)\cdots(n-r+1).
    \]
  \end{block}
  \begin{exampleblock}{Complete case}
    If $r=n$, then $P(n,n)=n!$. Order distinguishes every arrangement.
  \end{exampleblock}
\end{frame}

% 07 — Combinations
\begin{frame}{Combinations: when order does not matter}
  \small
  \begin{block}{Definition}
    The number of subsets of size $r$ chosen from $n$ objects is
    \[
      \binom{n}{r}=\frac{P(n,r)}{r!}=\frac{n!}{r!(n-r)!}.
    \]
  \end{block}
  \begin{exampleblock}{Key idea}
    Each unordered selection produces $r!$ ordered arrangements; division by
    $r!$ removes that multiplicity.
  \end{exampleblock}
\end{frame}

% 08 — Properties and decision
\begin{frame}{Properties and model selection}
  \footnotesize
  \begin{columns}[T]
    \begin{column}{0.48\textwidth}
      \begin{block}{Binomial coefficient}
        \[
          \binom{n}{r}=\binom{n}{n-r},\qquad
          \binom{n}{0}=\binom{n}{n}=1.
        \]
      \end{block}
    \end{column}
    \begin{column}{0.48\textwidth}
      \begin{alertblock}{Diagnostic question}
        \begin{itemize}\itemsep=0.05cm
          \item Does order change the outcome? $\to P(n,r)$.
          \item Does only the chosen set matter? $\to \binom{n}{r}$.
          \item Are there successive stages? $\to n_1\cdots n_k$.
        \end{itemize}
      \end{alertblock}
    \end{column}
  \end{columns}
\end{frame}

% 09 — Computational bridge 1/4
\begin{frame}{Computational bridge 1/4: exact arithmetic}
  \small
  \begin{block}{Reproducible verification}
    The formulas can be checked with integers, without rounding:
    \[
      8!=40320,\qquad 5!=120,\qquad 10!=3628800.
    \]
  \end{block}
  \begin{exampleblock}{Control principle}
    Compute factorials first and simplify ratios afterward. This separates a
    counting error from a numerical-approximation error.
  \end{exampleblock}
\end{frame}

% 10 — Computational bridge 2/4
\begin{frame}{Computational bridge 2/4: permutations and combinations}
  \footnotesize
  \begin{columns}[T]
    \begin{column}{0.48\textwidth}
      \begin{block}{Permutation}
        \[
          P(8,3)=8\cdot7\cdot6=336.
        \]
        Gold, silver, and bronze produce distinct ordered outcomes.
      \end{block}
    \end{column}
    \begin{column}{0.48\textwidth}
      \begin{block}{Combination}
        \[
          \binom{10}{4}=\frac{10!}{4!6!}=210.
        \]
        Swapping members does not create a new committee.
      \end{block}
    \end{column}
  \end{columns}
\end{frame}

% 11 — Computational bridge 3/4
\begin{frame}{Computational bridge 3/4: the sample space}
  \small
  \begin{block}{Five-card hands}
    For a hand of $5$ cards from a $52$-card deck,
    \[
      \lvert\Omega\rvert=\binom{52}{5}=2{,}598{,}960.
    \]
    The order in which cards arrive does not change the final hand.
  \end{block}
  \pause
  \begin{alertblock}{Favorable outcomes}
    One suit has $\binom{13}{5}$ hands and there are four suits, so
    $\lvert A\rvert=4\binom{13}{5}=5{,}148$.
  \end{alertblock}
\end{frame}

% 12 — Computational bridge 4/4
\begin{frame}{Computational bridge 4/4: numerical checks}
  \scriptsize
  \begin{center}
    \begin{tabular}{lrr}
      \hline
      Count & Exact result & Interpretation \\
      \hline
      $4\cdot6\cdot3$ & $72$ & menus \\
      $P(8,3)$ & $336$ & ordered medals \\
      $\binom{10}{4}$ & $210$ & committee without offices \\
      $4\binom{13}{5}$ & $5{,}148$ & same-suit hands \\
      \hline
    \end{tabular}
  \end{center}
  \vspace{0.1cm}
  Every value follows from the section formulas and can be reproduced with
  integer arithmetic independently of the deck language.
\end{frame}

% 13–14 — Example 1
\begin{frame}{Example 1/4 --- statement: restaurant menu}
  \small
  A restaurant offers $4$ appetizers, $6$ main courses, and $3$ desserts.
  If a diner chooses exactly one option from each category, how many different
  menus can be formed?
  \vfill
  \begin{alertblock}{Model}
    Three successive stages; every option combines with all later options.
  \end{alertblock}
\end{frame}

\begin{frame}{Example 1/4 --- solution}
  \small
  \begin{block}{Apply the multiplication principle}
    \[
      N=4\cdot6\cdot3=72.
    \]
  \end{block}
  \begin{exampleblock}{Interpretation}
    We do not list menus: each menu is a path through three decisions, so we
    count all possible paths.
  \end{exampleblock}
\end{frame}

% 15–16 — Example 2
\begin{frame}{Example 2/4 --- statement: medals}
  \small
  In a race with $8$ runners, in how many ways can gold, silver, and bronze be
  awarded, assuming there are no ties?
  \vfill
  \begin{alertblock}{Diagnosis}
    A medal assignment is ordered: exchanging two runners creates a different
    assignment.
  \end{alertblock}
\end{frame}

\begin{frame}{Example 2/4 --- solution}
  \small
  \begin{block}{Permutation of $8$ taken $3$ at a time}
    \[
      P(8,3)=\frac{8!}{(8-3)!}=8\cdot7\cdot6=336.
    \]
  \end{block}
  \begin{exampleblock}{Decision}
    Use $P(n,r)$ because gold, silver, and bronze are not interchangeable
    positions.
  \end{exampleblock}
\end{frame}

% 17–18 — Example 3
\begin{frame}{Example 3/4 --- statement: committee}
  \small
  A committee of $4$ people is formed from $10$ candidates. All positions are
  equivalent. How many different committees can be formed?
  \vfill
  \begin{alertblock}{Diagnosis}
    The committee is a subset: the order in which its members are selected
    does not change the outcome.
  \end{alertblock}
\end{frame}

\begin{frame}{Example 3/4 --- solution}
  \small
  \begin{block}{Combination of $10$ taken $4$ at a time}
    \[
      \binom{10}{4}=\frac{10!}{4!6!}=210.
    \]
  \end{block}
  \begin{exampleblock}{Interpretation}
    The $4!$ orderings of each committee describe the same committee and are
    removed by the combination formula.
  \end{exampleblock}
\end{frame}

% 19–20 — Example 4
\begin{frame}{Example 4/4 --- statement: same-suit hand}
  \small
  A five-card hand is drawn from a standard $52$-card deck. What is the
  probability that all five cards have the same suit?
  \vfill
  \begin{alertblock}{Model}
    Both the sample space and favorable outcomes use combinations because the
    order in which cards arrive does not matter.
  \end{alertblock}
\end{frame}

\begin{frame}{Example 4/4 --- solution}
  \footnotesize
  \begin{block}{Classical probability}
    \[
      P(A)=\frac{4\binom{13}{5}}{\binom{52}{5}}
      =\frac{5{,}148}{2{,}598{,}960}\approx0.00198.
    \]
  \end{block}
  \begin{exampleblock}{Reading the result}
    About one in every $505$ hands has the same suit; the denominator counts
    all hands and the numerator counts favorable hands.
  \end{exampleblock}
\end{frame}

% 21 — Synthesis
\begin{frame}{Section synthesis}
  \small
  \begin{itemize}\itemsep=0.12cm
    \item The product rule counts successive stages.
    \item Permutations model ordered outcomes.
    \item Combinations model unordered selections.
    \item Classical probability is the ratio of favorable outcomes to all
      equiprobable outcomes.
  \end{itemize}
\end{frame}

% 22 — Transition
\begin{frame}{Transition: from counting to conditional probability}
  \small
  \begin{block}{What we now know}
    We can construct finite sample spaces and measure events when all outcomes
    are equally likely.
  \end{block}
  \begin{alertblock}{Next section --- 02.04}
    Conditional probability restricts the sample space to the available
    information and prepares Bayes' rule.
  \end{alertblock}
\end{frame}

\end{document}
